\documentclass[10pt]{article}

\usepackage{amsmath}
\usepackage{amssymb}
\usepackage{amsthm}
\usepackage{ mathrsfs }
\usepackage{ mathtools}
\usepackage{enumerate}
\usepackage{bbm}
\usepackage{lipsum}
\usepackage{fancyhdr}
\usepackage{tikz-cd} 
\usepackage{adjustbox} 
\usetikzlibrary{arrows}
\newcommand{\midarrow}{\tikz \draw[-triangle 90] (0,0) -- +(.1,0);}
\tikzset{commutative diagrams/.cd,
mysymbol/.style={start anchor=center,end anchor=center,draw=none}
}
\newcommand\comsymb[2][\circlearrowleft]{%
  \arrow[mysymbol]{#2}[description]{#1}}

\newtheorem{theorem}{Theorem} [section] 
\newtheorem{proposition}{Proposition}[section] 
\newtheorem{definition}{Definition}[section] 
\newtheorem{lemma}{Lemma}[section] 
\newtheorem{notation}{Notation}[section] 
\newtheorem{remark}{Remark}[section] 
\newtheorem{corollary}{Corollary} [section] 
\newtheorem{terminology}{Terminology}[section] 
\newtheorem{fact}{Fact}[section] 
\newtheorem{example}{Example}[section] 
\newtheorem{claim}{Claim}
\newenvironment{claimproof}[1]{\par\noindent\underline{Proof:}\space#1}{\hfill $\blacksquare$}

\DeclareMathOperator{\tmod}{mod}
\DeclareMathOperator{\triv}{triv}
\DeclareMathOperator{\ind}{ind}
\DeclareMathOperator{\straw}{s.t.}
\DeclareMathOperator{\ev}{ev}
\DeclareMathOperator{\pr}{pr}
\DeclareMathOperator{\Aut}{Aut}
\DeclareMathOperator{\Map}{Map}
\DeclareMathOperator{\Stab}{Stab}
\DeclareMathOperator{\Ind}{Ind}
\DeclareMathOperator{\GL}{GL}
\DeclareMathOperator{\SL}{SL}
\DeclareMathOperator{\SO}{SO}
\DeclareMathOperator{\Sp}{Sp}
\DeclareMathOperator{\esssup}{esssup}
\DeclareMathOperator{\abs}{abs}
\DeclareMathOperator{\rel}{rel}
\DeclareMathOperator{\diam}{diam}
\DeclareMathOperator{\rank}{rank}
\DeclareMathOperator{\Hom}{Hom}
\DeclareMathOperator{\Homk}{Hom_{k-alg}}
\DeclareMathOperator{\Ker}{Ker}
\DeclareMathOperator{\Image}{Im}
\DeclareMathOperator{\Dom}{Dom}
\DeclareMathOperator{\grad}{grad}
\DeclareMathOperator{\divergence}{div}
\DeclareMathOperator{\rk}{rk}
\DeclareMathOperator{\Span}{Span}
\DeclareMathOperator{\mSpec}{mSpec}
\DeclareMathOperator{\MaxSpec}{MaxSpec}
\DeclareMathOperator{\interior}{int}
\DeclareMathOperator{\supp}{supp}
\DeclareMathOperator{\id}{id}
\DeclareMathOperator{\sgn}{sgn}
\DeclareMathOperator{\Mat}{Mat}
\DeclareMathOperator{\PD}{PD}
\DeclareMathOperator{\ext}{ext}
\DeclareMathOperator{\vol}{vol}
\DeclareMathOperator{\inte}{int}
\DeclareMathOperator{\diverg}{div}
\DeclareMathOperator{\curl}{curl}
\DeclareMathOperator{\image}{im}
\DeclareMathOperator{\dR}{dR}
\DeclareMathOperator{\Ob}{Ob}
\DeclareMathOperator{\Mnfd}{Mnfd}
\DeclareMathOperator{\OpEmb}{OpEmb}
\DeclareMathOperator{\Spec}{Spec}
\DeclareMathOperator{\Spa}{Spa}
\DeclareMathOperator{\Sper}{Sper}
\DeclareMathOperator{\op}{op}
\DeclareMathOperator{\con}{con}
%\newcommand*{\name}[\num_arguments][default values]{{\color{#1}\Large #2}}
\newcommand{\defeq}{\vcentcolon=}
\newcommand{\Ouv}{\mathcal{O}}
\newcommand{\norm}[1]{\Vert #1 \Vert}
\newcommand{\opNorm}[2]{\norm{#1}_{#2\to#2}}
\newcommand{\normL}[3]{\norm{#1}_{L^{#2}(#3)}}
\newcommand{\N}[0]{\mathbb{N}}
\newcommand{\m}[0]{\mathfrak{m}}
\newcommand{\R}[0]{\mathbb{R}}
\newcommand{\Z}[0]{\mathbb{Z}}
\newcommand{\C}[0]{\mathbb{C}}
\newcommand{\Q}[0]{\mathbb{Q}}
\newcommand{\Qc}[0]{\mathfrak{Q}_C}
\newcommand{\F}[0]{\mathbb{F}}
\newcommand{\G}[0]{\mathbb{G}}
\newcommand{\A}[0]{\mathbb{A}}
\newcommand{\T}[0]{\mathbb{T}}
\newcommand{\Para}[0]{\mathbb{P}}
\newcommand{\M}[0]{\mathcal{M}}
\newcommand{\maniN}[0]{\mathcal{N}}
\newcommand{\cinf}[0]{\mathcal{C}^\infty}
\newcommand{\torus}[0]{\mathbb{T}}
\newcommand{\sep}[0]{\:|\:}
\newcommand{\pd}[2]{{\frac {\partial #1} {\partial #2}}}
\newcommand{\compinc}[0]{\subset \subset}
\newcommand{\sH}[0]{\mathscr{H}}
\newcommand{\ab}[1]{\vert #1 \vert}
\newcommand{\st}[0]{\:\straw\:}
\newcommand{\g}[0]{\mathfrak{g}}
\newcommand{\p}[0]{\mathfrak{p}}
\newcommand{\U}[0]{\mathfrak{U}}
\newcommand{\fib}[1]{%
  \mathbin{\mathop{\times}\limits_{#1}}%
}
\newcommand{\tens}[1]{%
  \mathbin{\mathop{\otimes}\displaylimits_{#1}}%
}

\title{Rigid Analytic Geometry Week 3 Solution}
\author{So Murata}
\date{SoSe 26, University of Bonn}


\begin{document}
\maketitle

\par\textbf{Problem 1}\\
As we have seen in the previous sheet that for a $K$-norm $\norm{\cdot}$ on $L$, there is $C>0$ such that
\begin{equation*}
    \forall x,y\in L, \norm{xy}\leq C\norm{x}\norm{y}
\end{equation*}
Let $\norm{\cdot}'$ be another norm such that 
\begin{equation*}
    \norm{\cdot}' = C\norm{\cdot}.
\end{equation*}
Then using the equation,
\begin{equation*}
    {\frac 1 C}\norm{xy}'=\norm{xy}\leq C\norm{x}\norm{y} = {\frac C {C^2}}\norm{x}'\norm{y}'.
\end{equation*}
Thus,
\begin{equation*}
    \norm{xy}'\leq \norm{x}'\norm{y}'.
\end{equation*}
\par\textbf{Problem 2}\\
Suppose the limit,
\begin{equation*}
    \ab{l}_L = \lim_{n\to\infty}\sqrt[n]{\norm{l^n}_1}
\end{equation*}
exists for some norm $\norm{\cdot}_1$. Then for any norm $\norm{\cdot}_2$, we have,
\begin{equation*}
    c\norm{x}_1\leq \norm{x}_2\leq C\norm{x}_1,
\end{equation*}
for some $c,C>0$. Thus using the sandwich theorem, we obtain 
\begin{equation*}
    \lim_{n\to\infty}\sqrt[n]{\norm{l^n}_2}
\end{equation*}
also exists and equals to $\lim_{n\to\infty}\sqrt[n]{\norm{l^n}_1}$. 
Let $a_n = \log \norm{l^n}$. Then we have,
\begin{equation*}
    a_{m+n} \leq a_m+a_n.
\end{equation*}
Using Fekete's lemma, there is a limit $\lim_{n\to\infty}{\frac {a_n} n}$ and thus there is a limit for $\lim_{n\to\infty}\sqrt[n]{\norm{l^n}}$.
\par\textbf{Problem 3}
\begin{equation*}
    \ab{xy}_L=\lim_{n\to\infty}\sqrt[n]{\norm{(xy)^n}}\leq \lim_{n\to\infty}\sqrt[n]{\norm{x^n}\norm{y^n}} = \ab{x}_L\ab{y}_L.
\end{equation*}
Suppose $\norm{x} \leq \norm{y}$ and $\norm{y}\not=0$. Then, 
\begin{equation*}
    \ab{x}_L=\lim_{n\to\infty}\sqrt[n]{\norm{x^n}}\leq\lim_{n\to\infty}\sqrt[n]{\norm{y}} = \ab{y}_L.
\end{equation*}
\begin{align*}
    \ab{x+y}_L&=\lim_{n\to\infty}\sqrt[n]{\left\Vert y^n\left({\frac x y }+1\right)^n\right\Vert}\leq \lim_{n\to\infty}\sqrt[n]{\norm{y^n}} = \ab{y}_L,\\
\end{align*}
Thus the strong triangle equality holds. For any $\lambda\in K$, 
\begin{equation*}
    \ab{\lambda x} = \lim_{n\to\infty}\sqrt[n]{\lambda^n\norm{x^n}} = \ab{\lambda}_K\lim_{n\to\infty}\sqrt[n]{\norm{x^n}} = \ab{\lambda}_K\ab{x}_L.
\end{equation*}
Finally, $\ab{0}_L=0$ as $\norm{0^n}=0$.
\par\textbf{Problem 4}\\
This amounts to show that $\ab{1}_L = \ab{1}_K$. Take a supremum norm $\norm{\cdot}_\infty$, then $\norm{1}_\infty=1$. Since the choice of norms is independent when defining $\ab{\cdot}_L$, we see,
\begin{equation*}
    \ab{1}_L = \lim_{n\to\infty}\sqrt[n]{\norm{1}_{\infty}} = 1.
\end{equation*}
The general case follows tha $\ab{k}_L = \ab{k}_K\ab{1}_L$ for all $k\in K$.
\par\textbf{Problem 5}\\
From Problem 1, we have,
\begin{equation*}
    \norm{x}= \norm{l^{-1}lx}\leq \norm{m_{l^{-1}}}_{\op}\norm{lx}.
\end{equation*}
Multiply ${\frac 1 {\norm{m_{l^{-1}}}_{\op}^2\norm{l}}}$, we get,
\begin{equation*}
    \norm{m_{l^{-1}}}^{-1}{\frac {\norm{m_{l^{-1}}}^{-1}}{\norm{l}}}\norm{x}\leq {\frac {\norm{m_{l^{-1}}}^{-1}} {\norm{l}}}\norm{lx}.
\end{equation*}
Define $\norm{\cdot}' = {\frac {\norm{m_{l^{-1}}}^{-1}} {\norm{l}}}\norm{\cdot}$, then the above equation turns into,
\begin{equation*}
    \norm{l}'\norm{x}'\leq \norm{lx}.
\end{equation*}
Using the independence of choices of norms defining $\ab{\cdot}_L$, we have,
\begin{equation*}
    \ab{l}_L\ab{x}\leq \ab{lx}_L\leq \ab{l}_L\ab{x}_L.
\end{equation*}
Thus $\ab{\cdot}_L$ is such norm.
\par\textbf{Problem 6}\\
We have $\ab{1}_L = 1,\ab{0}_L = 0$. It is also multiplicative. As it is a $K$-norm, we see this is ultrametric.
\par\textbf{Problem 7}\\
Since $X$ is spectral $\mathfrak{Q}_c$ is a topology base of it. Note that for any $z\in S_x$, we have $x\in\Omega\in\Qc(X)$ then $z\in\Omega$. Thus,
\begin{equation*}
    S_x = \bigcap_{\substack{\Omega\in\Qc(X)\\ x\in\Omega}}\Omega.
\end{equation*}
From this 
\begin{equation*}
    S_x\cap S_y = \bigcap_{\substack{\Omega\in\Qc(X)\\ x\in\Omega}}\Omega\cap \bigcap_{\substack{\Theta\in\Qc(X)\\ y\in\Theta}}\Theta = \bigcap_{\substack{x\in\Omega\\y\in\Theta}}\Omega\cap\Theta.
\end{equation*}
\par\textbf{Problem 8}\\
Note that $X_{\con}$ is compact. By the assumption,
\begin{equation*}
    \bigcap_{\Omega,\Theta}\Omega\cap\Theta = \emptyset.
\end{equation*}
In $X_{\con}$ we get,
\begin{equation*}
    \bigcup (X\backslash(\Omega\cap\Theta)) = X_{\con}
\end{equation*}
Using the compactness of $X_{\con}$, for some $\Omega_i,\Theta_i,i=1,\cdots,n$,
\begin{equation*}
    \bigcap_{i=1}^n \Omega_i\cap\Theta_i = \emptyset.
\end{equation*}
In particular $\Omega = \bigcap_{i=1}^n \Omega_i,\Theta_{i=1}^n\Theta_i$ are disjoint sets separating $x,y$.
\par\textbf{Problem 9}\\
Since $H$ is a subgroup $1\in H$ and $0\not\in H$. Thus $w(1) = 1,w(0) = 0$. For the multiplicativeness, note that $H$ is a group thus $ x\in H,y\not\in H\Rightarrow xy\not\in H$ by the cancellation law. Thus $w$ is multiplicative. 
$w$ satisfies the strong triangle inequality on $H$ and since it is $0$ outside $H$, it holds globally on $\Gamma$.
\end{document}